\cotefile{157-2}
\scanpage{353}
\mspage{92}

en \deleted{\unsure{quelque|général}}, sous réserve de traiter les 4 cas : \unsure{irréductibles}.

A) Cas $\deleted{\unsure{T}}\tau=\Delta_{\varepsilon}^{2}$, $\tau'=\Delta_{\varepsilon}^{2}$ \inserted{ou $\Delta_{\varepsilon}$}. \deleted{\unsure{Alors la cat.}}

La catégorie $\mathrm{Hom}(\Delta_{\varepsilon}^{2},\Delta_{\varepsilon})$ a deux éléments, \deleted{les} morphismes $\Delta_{\varepsilon}^{2}\longrightarrow\Delta_{\varepsilon}$ correspondant aux \deleted{\unsure{deux}} flèches de transition de $\Delta_{\varepsilon}^{2}$ ($=\mathbb{\Delta}_{2}^{+}$ ou $\mathbb{\Delta}_{2}^{-}$). Dans le cas $\deleted{\unsure{2}}\varepsilon=+1$, ce sont les deux applications
\diagram{\begin{tikzcd}[ampersand replacement=\&,column sep=small,row sep=small]
0\arrow[r]\&1\arrow[r,no head]\&2\\
\unsure{\text{a.}}\&0\arrow[r]\&1
\end{tikzcd}}
de $\{0,1,2\}$ dans $\{0,1\}$, qui appliquent $0$ en $0$, $2$ en $1$ et $1$ soit en $\unsure{1|0}$, \deleted{\unsure{soit en}} \inserted{$(F_0)$} soit en $1$ $(F_1)$. On a évidemment une flèche $F_1\longrightarrow F_0$ (le voyageur pressé \unsure{arrive|arriver} plus vite \unsure{...}). \deleted{\unsure{Donc} $\underline{\mathrm{Hom}}$} \inserted{D'une façon} \unsure{pour} $\varepsilon=-1$
\[
\underline{\mathrm{Hom}}(\Delta_{\varepsilon}^{2},\Delta_{\varepsilon})\simeq\mathbb{\Delta}^{1},\qquad \text{catégorie 0-connexe.}
\]

La catégorie $\mathrm{Hom}(\Delta_{\varepsilon}^{2},\Delta_{\varepsilon}^{2})$ a \unsure{trois} éléments : \deleted{\unsure{resp.}} \unsure{s'identifiant|s'identifient} à l'une des parties non \unsure{vide|vides} de l'ens $\mathrm{Tr}(\tau)$ à deux éléments des flèches de transition $\underline u_1,\underline u_2$. de $\tau$ : il y a $F_{\underline u_1}=F_0$, $\underline F_{u_2}$, $F_{\underline u_1,\underline u_2}$. Si $\varepsilon=+1$, il s'agit des applications croissantes (\unsure{sans|sens} plus) de $\{0,1,2\}$ dans $\{0,1,2\}$, appliquant $0$ \inserted{\unsure{?}} en $0$, $2$ en $2$, et $1$ en ce qu'elle voudra. Donc
\diagram{$\begin{array}{ccl}
1\mapsto2&\text{correspond :}&F_0\\
1\mathrel{\unsure{\mapsto|-}}1&\rule{2cm}{.35pt}&F_1\\
1\mapsto0&\rule{2cm}{.35pt}&F_2
\end{array}$}
et on a bien sûr
\diagram{\begin{tikzcd}[ampersand replacement=\&,column sep=small]
F_2\arrow[r,no head]\& F_1\arrow[r,no head]\& F_0
\end{tikzcd}}
\[
\underline{\mathrm{Hom}}(\Delta_{\varepsilon}^{2},\underline{\Delta}_{\varepsilon}^{2})\simeq\mathbb{\Delta}^{2},
\]
catégorie 0-connexe.

Dans ces deux cas, deux éléments de $\underline{\mathrm{Hom}}(\tau,\tau')$ sont reliés par une homotopie \uline{simple}.

% archive sheet number: 352
